\section{Methods}\label{sec:methods}

% hand-written study overview (no new numbers; every value appears in the generated tables or result READMEs)
\begin{table}[t]
\centering
\small
\caption{Study overview. GPT-2 and Gemma use different SAE families, hook conventions and data, so their results are
reported side by side and never ranked against each other.}
\label{tab:overview}
\begin{adjustbox}{max width=\linewidth}
\begin{tabular}{lll}
\toprule
 & GPT-2 Small & Gemma-2-2B \\
\midrule
Layers and site & all 12 layers, \texttt{blocks.$L$.hook\_resid\_pre} & layers 5, 12, 19, \texttt{blocks.$L$.hook\_resid\_post} \\
SAEs & in-house ReLU, TopK, JumpReLU (sigmoid surrogate); & public Gemma Scope JumpReLU, width 16{,}384 \\
 & width 3{,}072, target L0 128, 3 seeds (108 SAEs); & (pinned IDs \texttt{average\_l0\_68/82/73}) \\
 & data-scaled R2 SAEs, width 12{,}288; public \texttt{gpt2-small-res-jb} & \\
Datasets & OpenWebText; Geometry of Truth (\texttt{cities}, \texttt{neg\_cities}); RepE-LM & The Pile (validation file) \\
Levels evaluated & reconstruction, information, behaviour, causal & reconstruction, behaviour \\
Strongest evidence & \evlabel{test} (reconstruction), \evlabel{confirmed test} (R2 scaling; R4 behaviour-targeted fine-tuning), & \evlabel{validation} \\
 & \evlabel{validation} (information, causal) & \\
\bottomrule
\end{tabular}
\end{adjustbox}
\end{table}


\subsection{Models and activation sites}\label{sec:models}

\paragraph{GPT-2 Small.} GPT-2 Small has 12 transformer blocks with a residual stream of width $d=768$. All GPT-2
SAEs, in-house and public, are applied to the activation entering block $L$, \texttt{blocks.$L$.hook\_resid\_pre},
for $L=0,\dots,11$; at $L=0$ this is the token plus position embedding. The beginning-of-sequence (BOS) position is
excluded from SAE training and from every splice. In addition to our own SAEs we evaluate the public
\texttt{gpt2-small-res-jb} release \citep{bloom2024gpt2}, loaded with SAELens \citep{saelens2024}, which has 24{,}576
latents and was trained on 300M OpenWebText tokens. It serves only as an external reference and is never counted as one
of our improvements.

\paragraph{Gemma-2-2B.} Gemma-2-2B has 26 layers and a residual stream of width 2{,}304. We use the public Gemma Scope
residual-stream SAEs \citep{gemmascope2024} at layers 5, 12 and 19, chosen before any run as representative lower,
middle and later depths. Each is a width-16k JumpReLU SAE applied to the output of the layer,
\texttt{blocks.$L$.hook\_resid\_post}, which is the site these SAEs were trained on. The SAE identifiers
(\texttt{layer\_\{5,12,19\}/width\_16k/average\_l0\_\{68,82,73\}}), the SAE repository revision, the model revision
and the tokenizer revision were pinned in the pre-registration so that a later change of a ``canonical'' alias could
not substitute a different SAE. The model runs in bf16 with eager attention; SAE encoding and decoding are computed
in fp32. The \texttt{resid\_pre} convention for GPT-2 and the \texttt{resid\_post} convention for Gemma follow the
respective SAE families; they are not a layer mapping, and no correspondence between GPT-2 and Gemma layers is
claimed.

\subsection{Datasets and splits}\label{sec:data}

\paragraph{OpenWebText (OWT).} We sample 1{,}000 documents from \texttt{Skylion007/openwebtext} (streaming shuffle
with seed 42, 100--20{,}000 characters, exact deduplication) and split them 800/100/100 by document. Each document is
truncated to 128 tokens, which gives 101{,}600 training, 12{,}700 validation and 12{,}700 test tokens excluding BOS. OWT
is used to train the base SAEs and to measure reconstruction and behavioural splicing. The data-scaled SAEs of
Section~\ref{sec:r2} use 100M further OWT tokens from a shard that excludes every document in this sample.

\paragraph{Geometry of Truth (GoT).} The \texttt{cities} and \texttt{neg\_cities} statement sets
\citep{marks2023geometry}, 1{,}496 statements each, form the truth task. Splits are grouped by city (1{,}046 / 224 /
226 per set), so no city appears in two splits. The training split is further divided into city groups A and B:
group A may be used by training objectives, and group B is reserved for fitting independent evaluation probes. The
numerical GoT sets have a surface shortcut (the second number alone predicts the label with AUROC 0.84) and are used
only as a diagnostic.

\paragraph{RepE-LM persona task.} Our persona prompts were originally derived from representation engineering
\citep{zou2023representation} (``Pretend you're an honest/untruthful person''). That construction had a length
shortcut: \emph{untruthful} is three tokens longer than \emph{honest}, and prompt length alone predicts the label with
AUROC 0.80. Results obtained with those prompts are kept for provenance only and are not used as evidence in this
paper. We therefore built RepE-LM: 506 statements (train 417, validation 89, and a separate held-out test split) times
20 prompt variants, in which the two prompts of a pair differ in exactly one token at a fixed position. Templates
T1--T4 swap one honest/dishonest word pair each (truthful/deceptive, trustworthy/dishonest, sincere/lying,
facts/lies), each with an honest-synonym control; T5--T6 are held-out wordings; N1--N2 swap neutral words (tall/short,
red/blue) as a generic one-word-change control. Seven pre-registered construction gates verify the task
(Appendix Table~\ref{tab:m1}); in particular every shortcut feature, including prompt length, has AUROC 0.50.

\paragraph{The Pile.} For Gemma we use the validation file of the uncopyrighted Pile \citep{gao2020pile}
(\texttt{monology/pile-uncopyrighted}, pinned revision and SHA-256). With a seeded permutation (seed 42) over document
indices, the first 512 documents with at least 32 tokens form the evaluation set and the next 64 form a disjoint
reference set used only to compute the mean-ablation vector. Each document is tokenized with BOS and truncated to its
first 1{,}024 tokens, the context length used for Gemma Scope. The Pile test file was never downloaded. The earlier
exploratory analysis of The Pile used a different file (the training split of a parquet mirror) and informed only the
truncation policy.

\subsection{SAE families}\label{sec:saes}

An SAE with $m$ latents maps $x\in\mathbb{R}^d$ to
\begin{equation}
  z=\sigma\big(W_e(x-b_d)+b_e\big)\in\mathbb{R}^m,\qquad \hat{x}=W_d z+b_d,\qquad e=x-\hat{x}.
\end{equation}
Our base GPT-2 SAEs use three activation functions: ReLU with an L1 penalty, TopK with $k=128$ \citep{gao2024scaling},
and a JumpReLU with a learned per-latent threshold whose gradient uses a sigmoid straight-through surrogate
($T=0.1$). The surrogate is simpler than the published JumpReLU estimator \citep{rajamanoharan2024jumping} and is
labelled as such throughout. Every base SAE has 3{,}072 latents ($4\times$ expansion) and is trained for 15{,}000 Adam
steps (learning rate $3\times10^{-4}$, batch 256) on the OWT training tokens, with decoder columns renormalised after
each step and a multiplicative controller that steers L0 towards 128. With seeds 42, 1042 and 2042 this gives
$12\times3\times3=108$ SAEs. Architectures are compared only when every run's L0 is within $\pm10\,\%$ of the target;
107 of 108 runs qualify, and the layer-0 JumpReLU run with seed 42 (L0 73.6) is excluded from rankings but kept.

The data-scaled SAEs of Section~\ref{sec:r2} (R2) are trained at layers 1, 8 and 11 with 12{,}288 latents and 100M
tokens each: TopK with an auxiliary dead-latent loss \citep{gao2024scaling}, BatchTopK \citep{bussmann2024batchtopk},
and ReLU with an L1 coefficient controlled to L0 = 128, plus a TopK width ablation at 3{,}072 latents and a
frozen-random-decoder TopK baseline \citep{korznikov2026random}. Gemma Scope SAEs are pretrained JumpReLU SAEs used as
released, with no fine-tuning.

\subsection{Four levels of faithfulness}\label{sec:levels}

\paragraph{Reconstruction faithfulness.} Explained variance is
$\mathrm{EV}=1-\sum\|x-\hat{x}\|^2/\sum\|x-\bar{x}\|^2$ and $\mathrm{NMSE}=1-\mathrm{EV}$, computed over the
evaluated tokens; L0 is the mean number of active latents per token. Because L0 depends on the input distribution, we
report it separately on OWT, GoT and RepE-LM where task comparisons are made.

\paragraph{Information faithfulness.} A reader is a classifier for the task label (true/false for GoT,
honest/dishonest for RepE-LM). The AUROC gap is $\mathrm{AUROC}(x)-\mathrm{AUROC}(\hat{x})$. A \emph{frozen reader} is
fitted on original activations and applied unchanged to reconstructions; it measures whether the original readout
still works. A \emph{re-trained reader} is fitted on reconstructions of the training statements; it measures whether
the information is still linearly recoverable. Two references separate SAE-specific loss from generic
perturbation: isotropic noise with the norm of the SAE error, and PCA reconstructions fitted on OWT training
activations. Evaluation probes are always fitted on original activations or on training-split reconstructions, never
on the data or parameters of a training objective.

\paragraph{Behavioural faithfulness.} Activation splicing replaces $x$ by $\hat{x}$ at every non-BOS position of one
layer and continues the forward pass. We report the increase in next-token cross-entropy ($\Delta$CE, nats per
prediction), the KL divergence from the original next-token distribution, and loss recovered relative to mean ablation,
$(\mathrm{CE}_{\text{mean-abl}}-\mathrm{CE}_{\text{SAE}})/(\mathrm{CE}_{\text{mean-abl}}-\mathrm{CE}_{\text{orig}})$. The
\emph{norm-matched random-error control} replaces the SAE error at each token by a random direction of the same norm;
the ratio of SAE-induced to random-induced $\Delta$CE asks whether the SAE error is more harmful than an arbitrary error
of the same size.

\paragraph{Causal faithfulness.} For a pair of honest (H) and dishonest (U) RepE-LM prompts, the model's
next-token distributions differ by $\mathrm{KL}(p_H\|p_U)$. Copying an activation from the honest run into the
dishonest run at one layer and position group restores a share of that divergence, which locates where the model
carries the effect (Section~\ref{sec:causal}). At a located site we decompose the activation difference into its
reconstructed and error parts and measure how much of the effect each part restores when patched.

These levels are related but not equivalent. Reconstruction is an average over all directions; information concerns
a few task-relevant directions; behaviour concerns everything downstream computation reads; and causal preservation
concerns one specific effect. A good result at one level is therefore never used as evidence for another.

\subsection{Evidence labels and experimental discipline}\label{sec:evidence}

Every result carries one of five evidence labels. \evlabel{test} marks a held-out test split evaluated under a
written protocol. \evlabel{confirmed test} marks a pre-registered comparison first observed on validation data and
then confirmed in a single pre-registered held-out evaluation (M5, Section~\ref{sec:r2}). \evlabel{validation} marks a
pre-registered analysis evaluated on validation data; validation data were also used to compare methods, so effect
sizes may be optimistic. \evlabel{exploratory} marks analyses that were not pre-registered or that are diagnostic.
\evlabel{pre-registered only} marks experiments whose protocol is fixed but which have not produced results. Within
validation experiments, pre-registered outcome classes such as ``preliminary improvement'', ``no improvement'', ``not
creditable'' (a guard failed) and ``regression'' describe the result of a comparison; they do not raise its evidence
level.

From the within-layer reanalysis (M0) onwards, every experiment wrote its hypotheses, comparisons, success criteria
and script hashes to a pre-registration file before running, and every script refuses to overwrite existing outputs.
Test splits were never used for training, tuning, checkpoint selection or choosing between methods. The GoT test split
was scored in earlier frozen-probe evaluations and is treated as previously exposed (Appendix~\ref{app:exposure}).

\paragraph{Statistics.} SAEs that share seeds are compared pairwise by seed. Fine-tuning comparisons report the mean
paired difference, its standard deviation and a 95\,\% $t$-interval with two degrees of freedom; the pre-registered
``consistent'' rule requires the improvement on all three seeds and a mean difference larger than twice its standard
deviation. Comparisons between SAE families use Welch intervals over seeds and require all nine seed pairs to agree.
Information gaps use statement-grouped paired bootstrap intervals, with Benjamini--Hochberg correction across SAEs
within a task. The within-layer regressions of M0 use layer and architecture fixed effects with a Freedman--Lane
permutation within layer (5{,}000 permutations) and Holm correction over four primary tests. Causal site tests use
Holm correction over 65 cells against 200 matched-norm random vectors per cell. Gemma results are single pretrained
SAEs per layer, so their intervals are 95\,\% document-bootstrap intervals (1{,}000 resamples), not seed variability.
Raw MSE is never averaged across layers. MSE improvement is $100\times(\text{baseline}-\text{improved})/\text{baseline}$.
We also report the design's minimum detectable effect (MDE $=2.80\times$SE at two-sided $\alpha=0.05$ and power 0.80):
between-family differences in the GoT re-trained gap smaller than about 0.05--0.09 AUROC cannot be detected, and one
SAE's causal preservation share can be resolved to 0.013--0.017 at layer 1, 0.036--0.069 at layer 11 and
0.054--0.107 at layer 8.
